\newpage
\begin{center}
  \textbf{\large 2. ТЕОРЕТИЧЕСКОЕ ОПИСАНИЕ}
\end{center}
\refstepcounter{chapter}
\addcontentsline{toc}{chapter}{2. ТЕОРЕТИЧЕСКОЕ ОПИСАНИЕ}

\section{Постановка задачи}

В данной работе рассматривается двумерная задача всплытия газового пузырька в
окружающей жидкости. Расчетная область задается прямоугольником
\[
    (x,y)\in[-0.5,0.5]\times[0,2],
    \qquad
    t\in[0,3].
\]
В начальный момент пузырек имеет круговую форму и расположен около нижней части
области. Центр пузырька находится в точке
\[
    (x_c,y_c)=(0,0.5).
\]
Начальный радиус пузырька обозначается через \(R\). Именно этот параметр
используется для различения двух рассматриваемых нейросетевых постановок. В
первой архитектуре радиус считается фиксированным и не передается в нейронную
сеть. Во второй архитектуре радиус является дополнительным входным параметром,
поэтому одна сеть должна описывать семейство задач с разными начальными
радиусами.

В задаче присутствуют две фазы: фаза~1 --- лёгкая (пузырёк), фаза~2 --- тяжёлая
(окружающая жидкость). Численные значения плотностей, вязкостей, поверхностного
натяжения и референсного радиуса приведены в таблице~\ref{tab:params};
ускорение свободного падения направлено вертикально вниз:
$\mathbf{g}=(0,{-}0{,}98)^T$.
Разность плотностей приводит к тому, что лёгкий пузырёк поднимается вверх под
действием силы Архимеда. В процессе движения форма пузырька меняется, а вокруг
него формируются нестационарные поля скорости и давления.

На верхней и нижней границах области задается условие прилипания:
\[
    u=0,\qquad v=0.
\]
На боковых границах используется периодическое условие для скорости и давления:
\[
    u(-0.5,y,t)=u(0.5,y,t),
    \qquad
    v(-0.5,y,t)=v(0.5,y,t),
\]
\[
    p(-0.5,y,t)=p(0.5,y,t).
\]
Так как в несжимаемой постановке давление определяется с точностью до
аддитивной константы, на верхней границе дополнительно фиксируется опорное
безразмерное давление
\[
    p_\infty^*=1.
\]

Цель нейросетевой модели состоит в восстановлении четырех полей:
\[
    u(x,y,t),\qquad v(x,y,t),\qquad p(x,y,t),\qquad \alpha(x,y,t),
\]
где \(u\) и \(v\) --- компоненты скорости, \(p\) --- давление, а \(\alpha\) ---
объемная доля легкой фазы. Значение \(\alpha=1\) соответствует области
пузырька, значение \(\alpha=0\) --- окружающей жидкости, а переходная область
между этими значениями описывает положение интерфейса.

Данные для обучения и проверки получаются с помощью CFD-расчетов в Basilisk.
В этой работе CFD-решение играет две роли. Во-первых, оно задает данные о
положении интерфейса, на которых обучается поле \(\alpha\). Во-вторых, оно
используется как эталон для последующего сравнения предсказаний PINN по полям
скорости, давления и объемной доли. Для параметрической постановки используются
несколько симуляций с разными начальными радиусами. Обучающие радиусы задаются
множеством
\[
    R_{\mathrm{train}}\in\{0.18,0.22,0.25,0.28,0.32\},
\]
а для проверки интерполяции по параметру используются радиусы
\[
    R_{\mathrm{test}}\in\{0.20,0.30\}.
\]
Тестовые значения находятся внутри диапазона обучающих радиусов, но не входят в
обучающую выборку. Поэтому такая проверка показывает, способна ли
параметрическая сеть восстановить решение для радиуса, которого она не видела
во время обучения.

Для описания двухфазности используется объемная доля легкой фазы
\(\alpha(x,y,t)\). Через нее задаются плотность и вязкость смеси:
\begin{align}
    \rho(\alpha) &= \rho_2 + (\rho_1-\rho_2)\alpha, \\
    \mu(\alpha) &= \mu_2 + (\mu_1-\mu_2)\alpha.
    \label{eq:chapter2_rho_mu}
\end{align}
Такая запись означает, что внутри пузырька используются параметры фазы \(1\), в
окружающей жидкости --- параметры фазы \(2\), а в области интерфейса параметры
меняются вместе с \(\alpha\).

Для безразмерной записи вводятся масштабы
\[
    \Rbase=0.25,\qquad \Uref=1,\qquad \rhoref=\rho_2=1000.
\]
Безразмерные переменные определяются как
\[
    x^*=\frac{x}{\Rbase},\qquad
    y^*=\frac{y}{\Rbase},\qquad
    t^*=\frac{t\Uref}{\Rbase},\qquad
    p^*=\frac{p}{\rhoref \Uref^2}.
\]
В параметрической постановке радиус также подается в сеть в безразмерном виде:
\[
    R^*=\frac{R}{\Rbase}.
\]
Для обучающих и тестовых радиусов это дает
\[
    R^*_{\mathrm{train}}\in\{0.72,0.88,1.00,1.12,1.28\},
    \qquad
    R^*_{\mathrm{test}}\in\{0.80,1.20\}.
\]
Далее для краткости знак \(*\) у безразмерных переменных может опускаться.

\begin{table}[H]
\centering
\caption{Физические параметры задачи. Фаза~1 --- пузырёк, фаза~2 --- жидкость.}
\label{tab:params}
\begin{tabular}{lll}
\hline
Параметр & Обозначение & Значение \\
\hline
Плотность пузырька   & $\rho_1$ & $100$ \\
Плотность жидкости   & $\rho_2$ & $1000$ \\
Вязкость пузырька    & $\mu_1$  & $1$ \\
Вязкость жидкости    & $\mu_2$  & $10$ \\
Поверхностное натяжение & $\sigma$ & $24.5$ \\
Ускорение свободного падения & $g$ & $0.98$ \\
Эталонный радиус     & $R_0$    & $0.25$ \\
\hline
\end{tabular}
\end{table}

По этим параметрам вычисляются характерные безразмерные числа. Число Рейнольдса
(отношение инерционных сил к вязким) определяет режим течения:
\begin{equation}
    \mathrm{Re} = \frac{\rho_2\,U_0\,(2R_0)}{\mu_2}
    = \frac{1000\cdot 1\cdot 0.5}{10} = 50.
\end{equation}
Число Этвёша (Bond) характеризует соотношение гравитации и поверхностного
натяжения:
\begin{equation}
    \mathrm{Eo} = \frac{(\rho_2-\rho_1)\,g\,(2R_0)^2}{\sigma}
    = \frac{900\cdot 0.98\cdot 0.25}{24.5} \approx 9.0.
\end{equation}
Число Мортона связывает все четыре параметра задачи и удобно для классификации
режима деформации пузырька:
\begin{equation}
    \mathrm{Mo} = \frac{g\,\mu_2^4\,(\rho_2-\rho_1)}{\rho_2^2\,\sigma^3}
    = \frac{0.98\cdot 10^4\cdot 900}{10^6\cdot (24.5)^3} \approx 6\cdot 10^{-4}.
\end{equation}
Эти значения соответствуют режиму умеренно деформированного пузырька:
число Рейнольдса достаточно велико для развития инерционного движения, число
Этвёша порядка~9 указывает на заметное влияние гравитации на форму
интерфейса, а малое число Мортона свидетельствует о том, что вязкость не
подавляет деформацию. Такой режим является физически нетривиальным и хорошо
подходит для проверки численных методов.

Условие несжимаемости записывается в виде
\begin{equation}
    \nabla^*\cdot\mathbf{u}^*=0.
    \label{eq:chapter2_mass}
\end{equation}
Поле объемной доли переносится потоком:
\begin{equation}
    \frac{\partial\alpha}{\partial t^*}
    +\mathbf{u}^*\cdot\nabla^*\alpha=0.
    \label{eq:chapter2_alpha_transport}
\end{equation}
Уравнение импульса с учетом вязкости, поверхностного натяжения и гравитации
имеет вид
\begin{align}
    \rho^*
    \left(
        \frac{\partial\mathbf{u}^*}{\partial t^*}
        +(\mathbf{u}^*\cdot\nabla^*)\mathbf{u}^*
    \right)
    &=
    -\nabla^*p^*
    +\nabla^*\cdot
    \left[
        C_\mu(\alpha)
        \left(\nabla^*\mathbf{u}^*+\nabla^*\mathbf{u}^{*T}\right)
    \right]
    \notag\\
    &\quad
    +C_\sigma\kappa\nabla^*\alpha
    +\rho^*\mathbf{C}_g.
    \label{eq:chapter2_momentum}
\end{align}
Здесь
\[
    \rho^*=\frac{\rho}{\rhoref},
    \qquad
    C_\mu(\alpha)=\frac{\mu(\alpha)}{\rhoref\Uref\Rbase},
    \qquad
    C_\sigma=\frac{\sigma}{\rhoref\Uref^2\Rbase},
\]
\[
    \mathbf{C}_g=\frac{\mathbf{g}\Rbase}{\Uref^2}.
\]
Кривизна интерфейса вычисляется по полю объемной доли:
\begin{equation}
    \kappa
    =
    -\nabla^*\cdot
    \frac{\nabla^*\alpha}{|\nabla^*\alpha|}.
    \label{eq:chapter2_curvature}
\end{equation}
В реализации при вычислении кривизны используется малое число
\(\varepsilon\), которое защищает выражение от деления на ноль в областях, где
градиент \(\alpha\) близок к нулю.

PINN обучается не только приближать данные, но и уменьшать невязки уравнений
\eqref{eq:chapter2_mass}--\eqref{eq:chapter2_momentum}. Невязка
несжимаемости имеет вид
\begin{equation}
    r_m=u_x+v_y.
\end{equation}
Невязка переноса объемной доли:
\begin{equation}
    r_\alpha=\alpha_t+u\alpha_x+v\alpha_y.
\end{equation}
Невязки импульса записываются покомпонентно. Обозначим
\[
    C_{\mu,x}=\frac{\mu_x}{\rhoref\Uref\Rbase},
    \qquad
    C_{\mu,y}=\frac{\mu_y}{\rhoref\Uref\Rbase},
    \qquad
    C_{g,y}=\frac{g_y\Rbase}{\Uref^2}.
\]
Тогда
\begin{align}
    r_x
    &=
    \rho^*
    \left(u_t+uu_x+vu_y\right)
    +p_x
    -C_\sigma\kappa\alpha_x
    -C_\mu\left(u_{xx}+u_{yy}\right)
    \notag\\
    &\quad
    -2C_{\mu,x}u_x
    -C_{\mu,y}\left(u_y+v_x\right),
    \label{eq:chapter2_rx}
    \\[0.4em]
    r_y
    &=
    \rho^*
    \left(v_t+uv_x+vv_y\right)
    +p_y
    -C_\sigma\kappa\alpha_y
    -C_\mu\left(v_{xx}+v_{yy}\right)
    \notag\\
    &\quad
    -\rho^*C_{g,y}
    -2C_{\mu,y}v_y
    -C_{\mu,x}\left(u_y+v_x\right).
    \label{eq:chapter2_ry}
\end{align}
Все производные в этих выражениях вычисляются средствами автоматического
дифференцирования в PyTorch. Это позволяет получать производные по входным
координатам непосредственно из нейронной сети, без построения отдельной
разностной сетки для PINN.

\section{Архитектура простой PINN}

Простая PINN рассматривается как базовая архитектура для восстановления одного
фиксированного режима всплытия пузырька. В этой постановке начальный радиус
\(R\) считается заданным заранее и не является входом нейронной сети. Поэтому
модель аппроксимирует одно конкретное решение уравнений двухфазного течения.

\subsection{Входы и выходы сети}

Для простой PINN нейронная сеть задает отображение
\[
    \mathcal{N}_{\theta}^{(0)}:
    (x^*,y^*,t^*)\mapsto (u^*,v^*,p^*,\alpha).
\]
На вход подаются две пространственные координаты и время. На выходе сеть
возвращает компоненты скорости, давление и объемную долю. Радиус пузырька в
такой архитектуре входит только в данные, на которых обучается модель, но не
передается в саму сеть. Поэтому после обучения модель связана с тем радиусом,
для которого были подготовлены обучающие точки.

\subsection{Структура нейронной сети}

В качестве базовой модели используется полносвязная нейронная сеть. Скрытый
блок состоит из нескольких линейных слоев с нелинейной функцией активации
\(\tanh\). Для простой PINN в экспериментальной части используется архитектура
с восемью скрытыми слоями по 350 нейронов в каждом:
\[
    (x^*,y^*,t^*) \rightarrow
    \underbrace{350 \rightarrow 350 \rightarrow \cdots \rightarrow 350}_{8\ \text{слоев}}
    \rightarrow
    (u^*,v^*,p^*,\alpha).
\]
Если обозначить скрытое представление последнего слоя через \(h\), то выходы
задаются отдельными головами:
\[
    u = W_u h + b_u,
    \qquad
    v = W_v h + b_v,
\]
\[
    p = \exp\left(\operatorname{clip}(W_p h+b_p,-20,20)\right),
    \qquad
    \alpha=\operatorname{sigmoid}(W_\alpha h+b_\alpha).
\]
Отдельные выходные головы позволяют сети независимо аппроксимировать физические
величины с разным масштабом и разной гладкостью. Сигмоида ограничивает
объемную долю диапазоном \(0\leq\alpha\leq1\), что согласуется с ее физическим
смыслом. Экспоненциальная параметризация давления обеспечивает положительность
предсказываемого давления в используемой постановке.

В программной реализации простая PINN может быть получена из параметрической
модели исключением радиуса из входного вектора. Тогда входной слой принимает
три величины вместо четырех, а остальные части архитектуры, физические невязки
и функция потерь остаются теми же.

\subsection{Функция потерь}

Функция потерь простой PINN объединяет данные, граничные условия и физические
невязки:
\begin{equation}
    \mathcal{L}^{(0)}
    =
    \mathcal{L}_{\alpha}
    +
    \mathcal{L}_{bc}
    +
    \lambda_m\mathcal{L}_{m}
    +
    \lambda_x\mathcal{L}_{mom,x}
    +
    \lambda_y\mathcal{L}_{mom,y}
    +
    \lambda_a\mathcal{L}_{\alpha,pde}.
    \label{eq:chapter2_loss_simple}
\end{equation}
Слагаемое \(\mathcal{L}_{\alpha}\) отвечает за совпадение предсказанной
объемной доли с данными CFD:
\begin{equation}
    \mathcal{L}_{\alpha}
    =
    \frac{1}{N_\alpha}
    \sum_{i=1}^{N_\alpha}
    \left(
        \alpha_\theta(x_i^*,y_i^*,t_i^*)-\alpha_i
    \right)^2.
\end{equation}
Граничное слагаемое \(\mathcal{L}_{bc}\) включает фиксацию давления на верхней
границе, периодичность боковых границ и условие нулевой скорости на твердых
границах. Физические слагаемые являются среднеквадратичными значениями
невязок:
\[
    \mathcal{L}_{m}
    =
    \frac{1}{N_f}\sum_{i=1}^{N_f}r_m(\mathbf{z}_i^f)^2,
    \qquad
    \mathcal{L}_{mom,x}
    =
    \frac{1}{N_f}\sum_{i=1}^{N_f}r_x(\mathbf{z}_i^f)^2,
\]
\[
    \mathcal{L}_{mom,y}
    =
    \frac{1}{N_f}\sum_{i=1}^{N_f}r_y(\mathbf{z}_i^f)^2,
    \qquad
    \mathcal{L}_{\alpha,pde}
    =
    \frac{1}{N_f}\sum_{i=1}^{N_f}r_\alpha(\mathbf{z}_i^f)^2.
\]
Веса физических невязок задают относительную важность разных уравнений. В
рассматриваемой реализации используются значения
\[
    (\lambda_m,\lambda_x,\lambda_y,\lambda_a)=(1,10,10,1).
\]

\subsection{Ограничения простой PINN}

Главное преимущество простой PINN состоит в том, что она решает одну задачу и
поэтому имеет более узкую область аппроксимации. Это упрощает обучение и делает
такую модель удобной базовой постановкой. Однако у нее есть принципиальное
ограничение: радиус пузырька не является параметром модели. Если требуется
получить решение для другого начального радиуса, необходимо заново подготовить
данные и обучить новую сеть. Следовательно, простая PINN подходит для
восстановления одного режима, но не решает задачу быстрого параметрического
моделирования.

\section{Архитектура параметрической PINN}

Параметрическая PINN строится как расширение простой архитектуры. Ее основная
идея состоит в том, чтобы обучить одну сеть не на одном фиксированном радиусе,
а на наборе CFD-симуляций с разными начальными радиусами. В результате сеть
должна аппроксимировать не отдельное решение, а зависимость решения от
параметра \(R\).

\subsection{Входы и выходы сети}

Параметрическая сеть задает отображение
\[
    \mathcal{N}_{\theta}^{(R)}:
    (x^*,y^*,t^*,R^*)\mapsto (u^*,v^*,p^*,\alpha).
\]
По сравнению с простой PINN входной вектор дополняется безразмерным радиусом
\(R^*\). Это позволяет использовать одну и ту же модель для разных начальных
размеров пузырька. Если простая PINN приближает функцию трех переменных, то
параметрическая PINN приближает функцию четырех переменных, где четвертая
переменная задает конкретную задачу из рассматриваемого семейства.

\subsection{Структура нейронной сети}

Структура параметрической сети совпадает с базовой MLP-идеей, но входной слой
принимает четыре величины. После входного слоя используется общий скрытый блок,
а затем четыре выходные головы для \(u\), \(v\), \(p\) и \(\alpha\). В текущем
проекте реализованы несколько конфигураций с функцией активации \(\tanh\):
\[
    (x^*,y^*,t^*,R^*) \rightarrow
    \underbrace{350 \rightarrow 350 \rightarrow \cdots \rightarrow 350}_{8\ \text{слоев}}
    \rightarrow
    (u^*,v^*,p^*,\alpha),
\]
\[
    (x^*,y^*,t^*,R^*) \rightarrow
    \underbrace{256 \rightarrow 256 \rightarrow \cdots \rightarrow 256}_{6\ \text{слоев}}
    \rightarrow
    (u^*,v^*,p^*,\alpha),
\]
\[
    (x^*,y^*,t^*,R^*) \rightarrow
    \underbrace{128 \rightarrow 128 \rightarrow 128 \rightarrow 128}_{4\ \text{слоя}}
    \rightarrow
    (u^*,v^*,p^*,\alpha).
\]
Такой набор конфигураций позволяет сравнивать влияние емкости сети на качество
восстановления полей. Более широкая и глубокая сеть обладает большей
выразительной способностью, но требует больше памяти и времени обучения.

\subsection{Функция потерь}

Функция потерь параметрической PINN имеет ту же структуру, что и
\eqref{eq:chapter2_loss_simple}. Отличие состоит в том, что каждая обучающая
точка содержит дополнительный параметр \(R^*\). Например, слагаемое по объемной
доле записывается как
\begin{equation}
    \mathcal{L}_{\alpha}^{(R)}
    =
    \frac{1}{N_\alpha}
    \sum_{i=1}^{N_\alpha}
    \left(
        \alpha_\theta(x_i^*,y_i^*,t_i^*,R_i^*)-\alpha_i
    \right)^2.
\end{equation}
Граничные условия также вычисляются в точках, дополненных соответствующим
радиусом. Для верхней границы используется условие
\[
    p_\theta(x_i^*,y_N^*,t_i^*,R_i^*)=p_\infty^*.
\]
Для боковых границ сравниваются предсказания в парных точках с одинаковыми
\(y^*\), \(t^*\) и \(R^*\):
\[
    u_\theta(\mathbf{z}_i^E)-u_\theta(\mathbf{z}_i^W)=0,
    \qquad
    v_\theta(\mathbf{z}_i^E)-v_\theta(\mathbf{z}_i^W)=0,
\]
\[
    p_\theta(\mathbf{z}_i^E)-p_\theta(\mathbf{z}_i^W)=0.
\]
Физические невязки \(r_m\), \(r_x\), \(r_y\) и \(r_\alpha\) считаются для
каждого значения радиуса из обучающего набора. Производные в невязках берутся
по координатам и времени; радиус при этом выступает как параметр задачи, а не
как дополнительная пространственная координата.

\subsection{Обучение на нескольких радиусах}

Для параметрической архитектуры обучающие данные объединяются из нескольких
CFD-файлов. Для каждого радиуса формируются множества точек:
\[
    \mathcal{D}_\alpha,\qquad
    \mathcal{D}_f,\qquad
    \mathcal{D}_{bc}.
\]
Множество \(\mathcal{D}_\alpha\) используется для обучения объемной доли,
\(\mathcal{D}_f\) --- для вычисления физических невязок, а
\(\mathcal{D}_{bc}\) --- для контроля граничных условий. В текущей реализации
для каждого выбранного временного слоя берутся точки около интерфейса и точки в
остальной области. Такое распределение точек важно, поскольку основные ошибки
в двухфазной задаче обычно сосредоточены рядом с границей фаз.

Из временной истории выбирается набор временных слоев, включающий ранние
моменты движения и последующую эволюцию пузырька. Ранние моменты важны, потому
что в начале расчета быстро формируются поля скорости и давления, а форма
пузырька начинает отклоняться от начальной круговой формы. На более поздних
этапах динамика становится более гладкой, поэтому временные слои могут быть
расположены реже.

Обучение выполняется алгоритмом Adam. Для основных экспериментов используются
две стадии:
\[
\begin{array}{c|cc}
    \hline
    \text{Стадия} & 1 & 2  \\
    \hline
    \text{Эпохи} & 5000 & 5000  \\
    \text{Шаг обучения} & 10^{-4} & 5\cdot10^{-5} \\
    \hline
\end{array}
\]
Данные перемешиваются в процессе обучения, а обучение выполняется мини-батчами
или наборами батчей, сформированными из разных типов точек. Такая схема
позволяет одновременно учитывать данные по интерфейсу, граничные условия и
физические ограничения.

\subsection{Преимущества и ограничения параметрической PINN}

Главное преимущество параметрической PINN состоит в том, что она не требует
обучать отдельную сеть для каждого радиуса. После обучения модель может
получать на вход новое значение \(R^*\) и выдавать приближенные поля
\(u^*\), \(v^*\), \(p^*\) и \(\alpha\). Это делает архитектуру более близкой к
практическому применению, где требуется быстро оценивать поведение системы при
изменении параметров.

Однако параметрическая постановка сложнее простой. Сеть должна одновременно
описать несколько режимов и согласовать между ними гладкую зависимость по
радиусу. Качество такой модели зависит от того, насколько хорошо обучающие
радиусы покрывают параметрическое пространство. Если новый радиус лежит далеко
от обучающего диапазона, ошибка может заметно возрасти. Поэтому в настоящей
работе проверяется именно интерполяция: тестовые радиусы не используются при
обучении, но находятся внутри диапазона обучающих значений.
